\section{Circuit Complexity --- Meyer Machine Setup}
%%% generated by the blueprint pipeline from TCSlib.Complexity.CircuitComplexity.MeyerMachineSetup
\subsection*{Circuit Complexity --- Meyer Machine Setup: Introduction}

Declarations of the Lean module \texttt{TCSlib.Complexity.CircuitComplexity.MeyerMachineSetup}.

\subsection*{Circuit Complexity --- Meyer Machine Setup: Results}

\begin{theorem}[Lookup after dropping a prefix]
\lean{Complexity.Meyer.getElem?_of_drop}
\leanok
\difficulty{1}
Let $y$ and $rest$ be Boolean words, and let $p$ and $i$ be natural numbers. If dropping
the first $p-1$ entries of $y$ gives $rest$, then the optional entry of $y$ at index
$p-1+i$ equals the optional entry of $rest$ at index $i$.
\end{theorem}

\begin{theorem}[Simulation of reading a field]
\lean{Complexity.Meyer.field_some}
\label{Complexity.Meyer.field_some}
\leanok
\uses{Complexity.Meyer.fpos, Complexity.Meyer.getElem?_of_drop, Complexity.Meyer.nbits, Complexity.Meyer.pcfg, Complexity.Meyer.pcfg_fA, Complexity.Meyer.pcfg_fB_false, Complexity.Meyer.pcfg_fB_true, Complexity.Meyer.readField, Complexity.Meyer.tabTM, Turing.MultiTapeTM.runFrom, Turing.MultiTapeTM.runFrom_succ_eq_step, Turing.MultiTapeTM.runFrom_zero}
\difficulty{6}
Let $M$ be a finite-state Boolean Turing machine, let $y$ be a Boolean word, and let $f$
and $bits$ be a Boolean and a bit vector of length equal to the number of selector types
of $M$. Suppose $rest,w,r,w_C,w_O$ are Boolean words, $p$ is natural with $1\le p$, the
suffix of $y$ from position $p-1$ is $rest$, and reading a field from $rest$ returns
$(w,r)$. After $2|w|+2$ steps of the tabulating machine for $M$, started in state
$fA(f,bits)$ at input position $\min(p,|y|+1)$ with auxiliary words $[],w_C,w_O$, its
entire configuration is the configuration at input position $\min(p+2|w|+2,|y|+1)$ with
state $cpy(bits)$ if $f$ is true and $fA(true,bits)$ otherwise, auxiliary words
$w_C,w_O$ changed to $w_C,w_O\mathbin{+\begin{bmatrix}w\end{bmatrix}}$ if $f$ is true,
and to $w_C\mathbin{+\begin{bmatrix}w\end{bmatrix}},w_O$ otherwise; the output remains
empty and the original machine's work tapes remain blank.
\end{theorem}

\begin{theorem}[Persistence of Rejection]
\lean{Complexity.Meyer.reject_persist}
\label{Complexity.Meyer.reject_persist}
\leanok
\uses{Complexity.Meyer.tabTM, Turing.Cfg, Turing.MultiTapeTM.runFrom, Turing.MultiTapeTM.runFrom_add, Turing.MultiTapeTM.runFrom_of_halt}
\difficulty{3}
Let $M$ be a finite Turing machine over the binary alphabet, let $y$ be a bit string,
and let $c$ be a configuration of the tabulating machine associated with $M$ on input
$y$. If after $s$ steps this machine is halted and its output is the one-bit word $0$,
then for every $s' \ge s$ it is still halted after $s'$ steps and its output remains
$0$.
\end{theorem}

\begin{theorem}[Halting on an invalid field]
\lean{Complexity.Meyer.field_none}
\label{Complexity.Meyer.field_none}
\leanok
\uses{Complexity.Meyer.fpos, Complexity.Meyer.getElem?_of_drop, Complexity.Meyer.nbits, Complexity.Meyer.pcfg, Complexity.Meyer.pcfg_fA, Complexity.Meyer.pcfg_fB_true, Complexity.Meyer.pcfg_reject, Complexity.Meyer.readField, Complexity.Meyer.tabTM, Turing.MultiTapeTM.runFrom, Turing.MultiTapeTM.runFrom_succ_eq_step, Turing.MultiTapeTM.runFrom_zero}
\difficulty{6}
Let $M$ be a finite-state Boolean Turing machine, let $y$ be a Boolean word, and let $f$
and $bits$ be a Boolean and a bit vector of length equal to the number of selector types
of $M$. For every Boolean word $rest$, natural number $p$, and Boolean words $w_C,w_O$,
suppose $1\le p$, $p+|rest|=|y|+1$, the suffix of $y$ beginning at position $p-1$ is
$rest$, and reading a field from $rest$ returns no result. Then, starting the tabulating
machine for $M$ in state $fA(f,bits)$ at input position $\min(p,|y|+1)$ with auxiliary
words $[],w_C,w_O$, there is a time $s\le |rest|+1$ at which it has halted and its
output is exactly $[\mathrm{false}]$.
\end{theorem}

\begin{theorem}[Copying a Word on the Simulated Tape]
\lean{Complexity.Meyer.cpy_run}
\label{Complexity.Meyer.cpy_run}
\leanok
\uses{Complexity.Meyer.acfg, Complexity.Meyer.bufferTape_write, Complexity.Meyer.fpos, Complexity.Meyer.getElem?_of_drop, Complexity.Meyer.moveInputPos_fpos, Complexity.Meyer.nbits, Complexity.Meyer.one_0_0, Complexity.Meyer.one_0_1, Complexity.Meyer.one_0_2, Complexity.Meyer.pcfg, Complexity.Meyer.pcfg_step, Complexity.Meyer.tabTM, Complexity.Meyer.tabTr, Complexity.Meyer.tact_idle_apply, Complexity.Meyer.wr_none, Turing.FinTM.bufferTape, Turing.MultiTapeTM.runFrom, Turing.MultiTapeTM.runFrom_succ_eq_step, Turing.MultiTapeTM.runFrom_zero, Turing.MultiTapeTM.step, Turing.moveInputPos_zero}
\difficulty{6}
Let $M$ be a finite Boolean Turing machine, let $y$ be a Boolean input word, and let
$bits$ be a Boolean assignment to the selector-bit positions of $M$. For any Boolean
word $rest$, natural number $p$, and Boolean words $w_X,w_C,w_O$, assume $1\le p$,
$p+|rest|=|y|+1$, and the suffix of $y$ beginning at index $p-1$ is $rest$. Starting the
machine that simulates $M$ in its copy state, with input head at the position obtained
by clamping $p$ to the valid input-head range and buffers containing $w_X,w_C,w_O$,
after $|rest|+1$ steps it is in the rewind-$X$ state at the right boundary, with the $X$
buffer containing $w_X$ followed by $rest$ and its head at position
$|w_X\mathbin{++}rest|-1$; the other buffers and their head positions are
unchanged, and the output is empty.
\end{theorem}

\begin{theorem}[First auxiliary read]
\lean{Complexity.Meyer.aux_read0}
\label{Complexity.Meyer.aux_read0}
\leanok
\difficulty{1}
Let $M$ be a finite-state Turing machine with $k$ work tapes, let $\alpha$ be a type,
let $f$ assign an $\alpha$-value to each work-tape index, and let $a,b,c$ be
$\alpha$-values. Extending $f$ by the three values $a,b,c$ on the three auxiliary
indices, the value at the first auxiliary index is $a$.
\end{theorem}

\begin{theorem}[Second auxiliary read]
\lean{Complexity.Meyer.aux_read1}
\label{Complexity.Meyer.aux_read1}
\leanok
\difficulty{1}
Let $M$ be a finite-state Turing machine with $k$ work tapes, let $\alpha$ be a type,
let $f$ assign an $\alpha$-value to each work-tape index, and let $a,b,c$ be
$\alpha$-values. Extending $f$ by the three values $a,b,c$ on the three auxiliary
indices, the value at the second auxiliary index is $b$.
\end{theorem}

\begin{theorem}[Third auxiliary read]
\lean{Complexity.Meyer.aux_read2}
\label{Complexity.Meyer.aux_read2}
\leanok
\difficulty{1}
Let $M$ be a finite-state Turing machine with $k$ work tapes, let $\alpha$ be a type,
let $f$ assign an $\alpha$-value to each work-tape index, and let $a,b,c$ be
$\alpha$-values. Extending $f$ by the three values $a,b,c$ on the three auxiliary
indices, the value at the third auxiliary index is $c$.
\end{theorem}

\begin{theorem}[Stored word at a predecessor index]
\lean{Complexity.Meyer.bufferTape_pred}
\label{Complexity.Meyer.bufferTape_pred}
\leanok
\uses{Turing.FinTM.bufferTape, Turing.FinTM.bufferTape_nat}
\difficulty{3}
Let $w$ be a Boolean word and let $j$ be a natural number with $j<|w|$. The tape storing
$w$ starting at cell zero contains $w_j$ at the integer cell $(j+1)-1$.
\end{theorem}

\begin{theorem}[Rewinding the First Auxiliary Tape]
\lean{Complexity.Meyer.rewX_run}
\label{Complexity.Meyer.rewX_run}
\leanok
\uses{Complexity.Meyer.acfg, Complexity.Meyer.acfg_step, Complexity.Meyer.aux_read0, Complexity.Meyer.bufferTape_pred, Complexity.Meyer.one_0_0, Complexity.Meyer.one_0_1, Complexity.Meyer.one_0_2, Complexity.Meyer.tabTM, Complexity.Meyer.tabTr, Complexity.Meyer.tact_idle_apply, Complexity.Meyer.wr_none, Turing.FinTM.bufferTape, Turing.FinTM.bufferTape_left, Turing.MultiTapeTM.runFrom, Turing.MultiTapeTM.runFrom_succ_eq_step, Turing.MultiTapeTM.runFrom_zero, Turing.moveInputPos_zero}
\difficulty{4}
Let $M$ be a finite Turing machine with Boolean symbols, let $y$ be a Boolean input
word, let $P$ be one of the $|y|+2$ input-head positions, and let $b$ be a Boolean
vector with one coordinate for each selector bit of $M$. Let $w$ be a Boolean word, let
$\tau_1,\tau_2:\mathbb{Z}\to\mathrm{Option}(\mathrm{Bool})$ be tape contents, and let
$z_1,z_2\in\mathbb{Z}$. For every natural number $j\le |w|$, after $j+1$ steps from the
configuration whose state is the first-tape-rewinding state, whose three auxiliary tapes
contain respectively $w$ starting at cell zero, $\tau_1$, and $\tau_2$, and whose
auxiliary head positions are $j-1,z_1,z_2$, the machine is in the next-tape-rewinding
state. The resulting tapes still contain those same contents, the input head remains at
$P$, the auxiliary head positions are $0,z_1-1,z_2$, and the output is empty.
\end{theorem}

\begin{theorem}[Rewinding the word tape]
\lean{Complexity.Meyer.rewC_run}
\label{Complexity.Meyer.rewC_run}
\leanok
\uses{Complexity.Meyer.acfg, Complexity.Meyer.acfg_step, Complexity.Meyer.aux_read1, Complexity.Meyer.bufferTape_pred, Complexity.Meyer.one_1_0, Complexity.Meyer.one_1_1, Complexity.Meyer.one_1_2, Complexity.Meyer.tabTM, Complexity.Meyer.tabTr, Complexity.Meyer.tact_idle_apply, Complexity.Meyer.wr_none, Turing.FinTM.bufferTape, Turing.FinTM.bufferTape_left, Turing.MultiTapeTM.runFrom, Turing.MultiTapeTM.runFrom_succ_eq_step, Turing.MultiTapeTM.runFrom_zero, Turing.moveInputPos_zero}
\difficulty{4}
Let $M$ be a finite Turing machine with Boolean symbols, let $y$ and $w$ be Boolean
words, let $P$ be an input-head position among the $|y|+2$ positions, and let $b$ be a
Boolean vector with one coordinate for each selector bit of $M$. Let
$	au_0,	au_2:\bbr	o	ext{Option}(	ext{Bool})$ be tape contents and let $z_0,z_2\in\bbr$.
For every natural number $j\le |w|$, after $j+1$ steps from the configuration in state
$\mathrm{rewC}$ with auxiliary tapes containing $\tau_0$, the word $w$ stored from cell
zero, and $\tau_2$, and head positions $z_0,j-1,z_2$, the machine is in state
$\mathrm{rewO}$ with the same tape contents, input-head position $P$, first and third
auxiliary head positions $z_0$ and $z_2-1$, and middle auxiliary head position $0$; the
output remains empty.
\end{theorem}

\begin{theorem}[Rewinding an auxiliary tape]
\lean{Complexity.Meyer.rewO_run}
\label{Complexity.Meyer.rewO_run}
\leanok
\uses{Complexity.Meyer.acfg, Complexity.Meyer.acfg_step, Complexity.Meyer.aux_read2, Complexity.Meyer.bufferTape_pred, Complexity.Meyer.one_2_0, Complexity.Meyer.one_2_1, Complexity.Meyer.one_2_2, Complexity.Meyer.tabTM, Complexity.Meyer.tabTr, Complexity.Meyer.tact_idle_apply, Complexity.Meyer.wr_none, Turing.FinTM.bufferTape, Turing.FinTM.bufferTape_left, Turing.MultiTapeTM.runFrom, Turing.MultiTapeTM.runFrom_succ_eq_step, Turing.MultiTapeTM.runFrom_zero, Turing.moveInputPos_zero}
\difficulty{4}
Let $M$ be a finite Turing machine with Boolean symbols, let $y$ be a Boolean input
word, let $P$ be an input-head position among the $|y|+2$ allowed positions, and let $b$
be a Boolean vector of length equal to the number of selector bits for $M$. Let $w$ be a
Boolean word, let $\tau_0,\tau_1:\mathbb{Z}\to\mathrm{Option}(\mathrm{Bool})$ be tape
contents, and let $z_0,z_1\in\mathbb{Z}$. For every $j\in\mathbb{N}$ with $j\le |w|$,
running the tabulating machine for $j+1$ steps from the configuration whose state is the
output-rewind state with vector $b$, whose input-head position is $P$, whose first two
auxiliary tapes and head positions are $(\tau_0,z_0)$ and $(\tau_1,z_1)$, and whose
third auxiliary tape stores $w$ from cell zero with head at $j-1$, yields the
configuration with state the decode state carrying $b$, the initial state of $M$, an
empty output register, and tag true. The input-head position, all tape contents, the
first two auxiliary head positions, and the empty output are unchanged; the third
auxiliary head is at cell zero.
\end{theorem}

\begin{theorem}[Equality of blank tapes]
\lean{Complexity.Meyer.bl_eq}
\label{Complexity.Meyer.bl_eq}
\leanok
\uses{Complexity.Meyer.bl, Turing.FinTM.bufferTape}
\difficulty{1}
The tape that is blank at every integer cell equals the tape obtained by storing the
empty word starting at cell zero.
\end{theorem}

\begin{theorem}[Empty Buffer Configuration]
\lean{Complexity.Meyer.pcfg_nil}
\label{Complexity.Meyer.pcfg_nil}
\leanok
\uses{Complexity.Meyer.St, Complexity.Meyer.acfg, Complexity.Meyer.bl, Complexity.Meyer.pcfg}
\difficulty{1}
Let $M$ be a finite Turing machine over Booleans, let $y$ be a Boolean word, and let
$st$ be a state of the constructed machine and $p$ a valid input-head position. The
configuration with all three auxiliary buffers empty and their heads at position zero is
equal to the configuration with all three tapes blank and their heads at position zero,
with the same state, input-head position, and empty output.
\end{theorem}

\begin{theorem}[Tabulation Machine Setup]
\lean{Complexity.Meyer.tabTM_setup}
\label{Complexity.Meyer.tabTM_setup}
\leanok
\uses{Complexity.Meyer.acfg, Complexity.Meyer.cpy_run, Complexity.Meyer.fcode, Complexity.Meyer.field_none, Complexity.Meyer.field_some, Complexity.Meyer.fpos, Complexity.Meyer.nbits, Complexity.Meyer.parseQ, Complexity.Meyer.pcfg, Complexity.Meyer.pcfg_nil, Complexity.Meyer.readField, Complexity.Meyer.readField_eq_some, Complexity.Meyer.rewC_run, Complexity.Meyer.rewO_run, Complexity.Meyer.rewX_run, Complexity.Meyer.sel_done, Complexity.Meyer.sel_reject, Complexity.Meyer.tabTM, Turing.FinTM.bufferTape, Turing.MultiTapeTM.initCfg, Turing.MultiTapeTM.runFrom, Turing.MultiTapeTM.runFrom_add}
\difficulty{6}
For a Boolean finite Turing machine $M$ and Boolean word $y$, if parsing $y$ as a query
with selector width determined by $M$ fails, the tabulation machine halts within $|y|+1$
steps and outputs the single bit $0$. If parsing succeeds with fields $bits,t,a,x$, then
for some number of steps at most $2|y|+4$, the machine reaches a configuration in its
decoding state, with the selector bits stored in its finite control, its input head at
the right boundary of $y$, its first three auxiliary work tapes holding the respective
words $x,t,a$ from cell zero, their heads at cell zero, its original work tapes blank
with heads at zero, and its output empty.
\end{theorem}
